\subsection{Groups generated by a compact subset}

Lemma 1. --- Let H be a locally compact group, and let R be one of the groups R or Z. Let \(\varphi\) be a continuous morphism from R to H. If \(\varphi\) is not a topological isomorphism of R onto a subgroup of H, then the image of R in H is relatively compact.

Let I be the image of \(\varphi\) . Replacing H by \(\bar{I}\) , we may assume that I is dense in H. We must then show that H is compact if \(\varphi\) is not a topological isomorphism of R onto I.

Suppose there exists a neighborhood V of \(e\) in H and an integer \(M > 0\) such that, for all \(t > \mathrm{M}\) in R, we have \(\varphi(t) \notin \mathrm{V}\). Then \(\varphi\) is injective: if \(\varphi(u) = e\), one has \(\varphi(nu) = e\) for every integer \(n \geqslant 1\), hence the set \(\mathbf{N}u\) is bounded in R, which means that \(u = 0\). The restriction of \(\varphi\) to \([-M, M] \cap R\) is therefore a homeomorphism onto its image, which contains \(V \cap I\). The restriction of \(\varphi^{-1}\) to \(V \cap I\) being continuous, it follows that \(\varphi\) is a topological isomorphism of R onto I.

Suppose now that \(\varphi\) is not a topological isomorphism from R onto I. Let W be a relatively compact open neighborhood of \(e\) in H, and V a symmetric neighborhood of \(e\) such that \(\mathrm{V}^2 \subset \mathrm{W}\). For every \(x \in \mathrm{H} = \overline{\mathrm{I}}\), there exists an element \(s \in \mathrm{R}\) such that \(x \in \varphi(s)\mathrm{V}\). From the preceding paragraph and the hypothesis on \(\varphi\), there exists \(t \in \mathrm{R}\) such that \(t > |s|\) and \(\varphi(t) \in \mathrm{V}\). We then have \(x \in \varphi(t + s)\varphi(t)^{-1}\mathrm{V} \subset \varphi(t + s)\mathrm{W}\), and \(t + s > 0\). Consequently, the open sets \(\varphi(u)\mathrm{W}\) for \(u > 0\) form an open covering of H. Since W is relatively compact, there exists an integer \(n \geqslant 1\) and elements \(u_1, \ldots, u_n\) of R, strictly positive, such that \(\overline{\mathrm{W}} \subset \bigcup_{1 \leqslant i \leqslant n} \varphi(u_i)\mathrm{W}\). Let U be the largest of the \(u_i\) .

Let \(x \in \mathrm{H}\) and let \(s = \inf \{t \in \mathrm{R} | t \geqslant 0, \varphi(t)x^{-1} \in \overline{\mathrm{W}}\}\) . Since \(\overline{\mathrm{W}}\) is compact, we then have \(\varphi(s)x^{-1} \in \overline{\mathrm{W}}\) . There exists an integer \(i\) such that \(\varphi(s)x^{-1} \in \varphi(u_i)\mathrm{W}\) , whence \(\varphi(s - u_i)x^{-1} \in \overline{\mathrm{W}}\) . The definition of \(s\) implies \(s - u_i < 0\) , whence \(s \leqslant \mathrm{U}\). It follows that \(\mathrm{H} = \varphi([0,\mathrm{U}] \cap \mathrm{R})\overline{\mathrm{W}}\) is compact.

Lemma 2. --- If G is generated by a compact subset V, then there exists an integer \(n \geqslant 0\) and a discrete subgroup D of G isomorphic to \(Z^{n}\) such that G/D is compact.

Up to replacing V by \(V \cup V^{-1}\) , we may assume that V is symmetric; the hypothesis then means that G is the union of the sets \(V^{n}\) where \(n \in N\) .

Since \(V^{2}\) is compact, there exists an integer \(k \geqslant 1\) and elements \(x_{1}, \ldots, x_{k} \in G\) such that \(V^{2} \subset \bigcup_{1 \leqslant i \leqslant k} x_{i}V\) . Let \(D_{0}\) be the subgroup of G generated by the family \((x_{i})_{1 \leqslant i \leqslant k}\). We have \(V^{2} \subset D_{0}V\), whence by induction \(V^{n} \subset D_{0}V\) for every integer \(n \geqslant 1\), and hence \(G = D_{0}V\) since V generates G. Let then J be a subset of \(\{1, 2, \ldots, k\}\) such that the subgroup D generated by the family \((x_{i})_{i \in J}\) is topologically isomorphic to \(\mathbf{Z}^{\mathrm{Card(J)}}\) and maximal for this property. Let us show that G/D is compact.

Let \(p\) be the canonical surjection from G onto G/D. Let \(i \in \{1, 2, \ldots, k\} - J\) . If the subgroup \(H_i\) of G/D generated by \(p(x_i)\) is topologically isomorphic to \(Z\) , the subgroup of G generated by D and \(x_i\) is discrete and the map \((d, n) \mapsto dx_i^n\) is an isomorphism from D × \(Z\) onto this subgroup, contrary to the maximality of J. Lemma 1 therefore implies that \(\overline{H}_i\) is compact. Hence G/D = \((\prod_{i \not\in J} \overline{H}_i)p(V)\) is compact.

Lemma 3. — Let A and B be commutative groups such that A is divisible. Let C be a subgroup of B and \(\varphi\) a morphism from C to A. There exists a morphism from B to A extending \(\varphi\) .

Let \(\mathcal{O}\) be the set of pairs \((\mathrm{X},f)\), where \(\mathrm{X}\) is a subgroup of B containing C and \(f\) a morphism from \(\mathrm{X}\) to A extending \(\varphi\). Order \(\mathcal{O}\) by the relation " \(\mathrm{X} \subset \mathrm{X}'\) and \(f'\) extends \(f\) ". One verifies that \(\mathcal{O}\) is inductive. Let \((\mathrm{X},f)\) be a maximal element of \(\mathcal{O}\) (E, III, p.~20, th. 2). If \(\mathrm{X} \neq \mathrm{B}\) , take an element \(b\) of B-X and let \(\mathrm{X}'\) be the subgroup generated by \(\mathrm{X}\) and \(b\) . The commutativity of B shows that \(\mathrm{X}'\) is the set of elements \(b^n x\) for \(n \in \mathbf{Z}\) and \(x \in \mathrm{X}\) . Suppose first that \(b^n \notin \mathrm{X}\) for every integer \(n \neq 0\) and define \(f'\) as a morphism from \(\mathrm{X}'\) to A by taking an arbitrary element \(y \in \mathrm{A}\) and setting \(f'(b^n x) = y^n f(x)\) for all \(n \in \mathbf{Z}\) and every \(x \in \mathrm{X}\) . Since A is commutative, \(f'\) is a morphism, and it extends \(f\) . Suppose now that there exists \(n \neq 0\) such that \(b^n \in \mathrm{X}\) and let \(m > 0\) such that \(m\mathbf{Z} = \{n \in \mathbf{Z} \mid b^n \in \mathrm{X}\}\) . Since A is divisible, there exists an element \(y \in \mathrm{A}\) such that \(y^m = f(b^m)\) . One then extends \(f\) to a morphism from \(\mathrm{X}'\) to A by \(f'(b^n x) = y^n f(x)\) for \(n \in \{0,1,\ldots,m-1\}\) and \(x \in \mathrm{X}\) . In both cases, \((\mathrm{X},f)\) would not be maximal. Hence we have \(\mathrm{X} = \mathrm{B}\) and the lemma is proved.

Remark. --- In the language of categories, the lemma says that divisible groups are injective objects in the category of commutative groups; cf.~A, VII, p.~53, exerc. 3.

PROPOSITION 1. --- The following conditions are equivalent:

\begin{enumerate}
\def\labelenumi{(\roman{enumi})}
\item
  G is generated by a compact subset;
\item
  there exist positive integers p and q and a compact group K such that G is isomorphic to \(R^{p} \times Z^{q} \times K\) ;
\item
  there exists an integer \(n \geqslant 0\) such that \(\widehat{G}\) is locally isomorphic to \(R^{n}\) ;
\item
  there exist positive integers p and q and a discrete group D such that \(\widehat{G}\) is isomorphic to \(R^{p} \times T^{q} \times D\) .
\item
  \(\Longrightarrow\) (iii): if G has property (i), there exists an integer \(n \geqslant 0\) and a subgroup D of G isomorphic to \(Z^{n}\) such that G/D is compact (lemma 2). Then \(D^{\perp}\) , which is identified with the dual of G/D, is discrete (th. 4 of II, p.~226 and prop. 18 of II, p.~233). Hence \(\widehat{G}\) is locally isomorphic to \(\widehat{G}/D^{\perp}\) , that is to say to \(\widehat{D}\) , which is isomorphic to \(T^{n}\) (th. 4 of II, p.~226 and prop. 18 of II, p.~233). Now \(T^{n}\) is locally isomorphic to \(R^{n}\) .
\item
  \(\Longrightarrow\) (iv): if \(\widehat{\mathbf{G}}\) is locally isomorphic to \(\mathbf{R}^n\) , there exists an integer \(p\) such that \(0 \leqslant p \leqslant n\) so that the neutral component \(\widehat{\mathbf{G}}_0\) of \(\widehat{\mathbf{G}}\) is an open subgroup isomorphic to \(\mathbf{R}^p \times \mathbf{T}^{n-p}\) (TG, VII, p.~13, th. 1). In particular, \(\widehat{\mathbf{G}}_0\) is a divisible group. Applying lemma 3 to the identity map of the subgroup \(\widehat{\mathbf{G}}_0\) of the group \(\widehat{\mathbf{G}}\) into the divisible group \(\widehat{\mathbf{G}}_0\) , there exists a morphism \(\pi\) of \(\widehat{\mathbf{G}}\) into \(\widehat{\mathbf{G}}_0\) which is the identity map on \(\widehat{\mathbf{G}}_0\) . Consequently, we have \(\pi \circ \pi = \pi\) , and \(\pi\) is a projector. It is continuous, since its restriction to the open subgroup \(\widehat{\mathbf{G}}_0\) is. Hence, \(\widehat{\mathbf{G}}\) is the direct product of \(\widehat{\mathbf{G}}_0\) and the subgroup \(\overline{\pi}^{-1}(e)\) , which is discrete since it is isomorphic to \(\widehat{\mathbf{G}} / \widehat{\mathbf{G}}_0\) (TG, III, p.~47, cor.) (iv) \(\Longrightarrow\) (ii): follows from prop. 2 of II, p.~206, from prop. 18 of II, p.~233 and from corollary 3 of II, p.~236.
\item
  \(\Longrightarrow\) (i): for every compact group K, the group \(R^{p} \times Z^{q} \times K\) is generated by the compact set \([0,1]^{p} \times \{0,1\}^{q} \times K\) .
\end{enumerate}

COROLLARY 1. --- Suppose that G is generated by a compact neighborhood of e.

\begin{enumerate}
\def\labelenumi{\alph{enumi})}
\item
  There exists a compact subgroup K of G and positive integers p and q such that G is isomorphic to \(R^{p} \times Z^{q} \times K\) .
\item
  Conversely, let K be a compact group, let p and q be positive integers, and let G be a group isomorphic to \(R^{p} \times Z^{q} \times K\). Then K is the unique maximal compact subgroup of G, and the integers \((p,q)\) are uniquely determined by G.
\end{enumerate}

Assertion \(a\) ) follows from prop. 1. Let then K be a compact group, and \(p\) , \(q\) positive integers. Suppose that G is isomorphic to the group \(\mathbf{R}^{p} \times \mathbf{Z}^{q} \times \mathrm{K}\) , and we identify G with this group. Under the canonical projection of G onto \(\mathbf{R}^{p} \times \mathbf{Z}^{q}\) , the image of every compact subgroup of G is a compact subgroup of \(\mathbf{R}^{p} \times \mathbf{Z}^{q}\); therefore it is reduced to the identity element. Therefore \(\mathrm{K}' \subset \mathrm{K}\) and K is the largest compact subgroup of G. The subgroup \(\mathbf{R}^{p} \times \mathrm{K}\) is also unique because \(\mathbf{R}^{p}\) is the neutral component of G/K. Taking into account TG, VII, p.~13, cor. 3, the integer \(p\) is uniquely determined by G. Since \(\mathrm{G}/(\mathbf{R}^{p} \times \mathrm{K})\) is isomorphic to \(\mathbf{Z}^{q}\) , the integer \(q\) is also uniquely determined by G.

Note. --- By duality, \(\hat{G}\) is isomorphic to \(R^{p} \times T^{q} \times D\) (prop. 3 (iv)), where the subgroups \(R^{p} \times T^{q}\) and \(T^{q}\), and the integers p and q, are uniquely determined.

COROLLARY 2. --- The following conditions are equivalent:

\begin{enumerate}
\def\labelenumi{(\roman{enumi})}
\item
  The groups G and \(\hat{G}\) are generated by compact subsets;
\item
  There exist positive integers n and m such that G is locally isomorphic to \(R^{m}\) and \(\hat{G}\) to \(R^{n}\) ;
\item
  There exist positive integers p, q, and r and a finite group A such that G is isomorphic to \(R^{p} \times T^{q} \times Z^{r} \times A\) ;
\item
  There exist positive integers p, q, and r and a finite group A such that \(\hat{G}\) is isomorphic to a product \(R^{p} \times Z^{q} \times T^{r} \times A\) .
\end{enumerate}

We have (i) \(\Leftrightarrow\) (ii) by Prop. 1, and (iii) \(\Leftrightarrow\) (iv) by duality, hence (iii) \(\Rightarrow\) (i). Finally, if (i) is true, then \(\widehat{\mathbf{G}}\) is isomorphic to \(\mathbf{R}^{p} \times \mathbf{T}^{r} \times \mathbf{D}\), where D is discrete (prop. 1). The group D is generated by a compact subset and is therefore finite. Consequently, there exists \(q \geqslant 0\) such that D is isomorphic to \(\mathbf{Z}^{q} \times \mathbf{A}\) , where A is a finite group (A, VII, p.~22, th. 3).

Note. --- With the notations of Cor. 2, if one identifies G with \(R^{p} \times T^{q} \times Z^{r} \times A\) the subgroup \(R^{p} \times T^{q}\) is the identity component of G, the subgroup \(T^{q} \times A\) is its largest compact subgroup and \(T^{q}\) is the identity component of it; the integers p, q, r are uniquely determined by G according to the preceding remark, and the group A is determined by G up to isomorphism.

PROPOSITION 2. --- Suppose G is compact. There exists a decreasing filtered family \((\mathrm{H}_{i})_{i\in\mathrm{I}}\) of closed subgroups of G such that

\begin{enumerate}
\def\labelenumi{\alph{enumi})}
\item
  the group G is identified with the projective limit of the \(\mathrm{G / H_i}\) ;
\item
  for every i, there exists an integer \(q \geqslant 0\) and a finite group A such that \(G/H_{i}\) is isomorphic to \(T^{q} \times A\) .
\end{enumerate}

Indeed, \(\widehat{\mathbf{G}}\) is discrete (prop. 18 of II, p.~233), hence is the union of an increasing filtered family \((\mathrm{D}_i)_{i\in \mathrm{I}}\) of finitely generated subgroups. Set \(\mathrm{H}_i = \mathrm{D}_i^\perp\); the group G is identified with the projective limit of the \(\mathrm{G / H_i}\) (II, p. 234, cor. 3), and \((\mathrm{H}_i)\) is a decreasing filtered family.

Let \(i \in I\) . There exists an integer \(q \geqslant 0\) and a finite group A such that the group \(D_{i}\) is isomorphic to \(Z^{q} \times A\) (A, VII, p.~22, th. 3), hence \(G/H_{i}\) is isomorphic to \(T^{q} \times \widehat{A}\) (cf.~corollary 1 of II, p.~232).

COROLLARY. --- If the group G is generated by a compact subset, then it is a projective limit of groups isomorphic to groups of the form \(R^{p} \times T^{q} \times Z^{r} \times A\) , where A is a finite group and p, q, r are positive integers.

By (ii) of prop. 1 of II, p.~246, the corollary follows from prop. 2.

